\documentclass[11pt,a4paper]{article}
\usepackage[margin=23mm]{geometry}
\usepackage{amsmath,amssymb,amsthm,booktabs,microtype,needspace}
\usepackage[T1]{fontenc}
\usepackage{lmodern}
\usepackage[colorlinks=true,linkcolor=blue,urlcolor=blue]{hyperref}
\hypersetup{pdftitle={Subset sums: joint covering losses and the audited bound 95},pdfauthor={Sum-of-Subset research project}}
\newcommand{\FS}{\operatorname{FS}}
\newcommand{\T}{\mathcal T}
\newcommand{\Z}{\mathbb Z}
\newtheorem{theorem}{Theorem}[section]
\newtheorem{lemma}[theorem]{Lemma}
\newtheorem{proposition}[theorem]{Proposition}
\title{Subset Sums in a Fourfold Sumset\\\large Joint covering losses and the audited bound $N\le95$}
\author{Sum-of-Subset research project}
\date{16 October 2026}
\begin{document}
\maketitle
\begin{abstract}
We exclude cardinality 96 and extend the independently audited chain to
$71\le N\le95$. A sharp four-extreme estimate reduces two coupled
budgets by two. The covering is then organized by actual union deficits:
leading twofold translates and retained onefold translates share one
loss budget. This handles three empty units and multiple large-loss
centers without a global ``at most one large loss'' assumption.
Positive low holes are transported to 14; 1,241 retained oriented
high-hole states are compressed into 26 endpoint groups.
All endpoint certificates are newly generated and checked by exact
integer arithmetic. The stronger archived bound 84 is preserved, but
its entire proof chain has not been independently replayed; this result
advances the audited chain and does not surpass 84.
\end{abstract}

\section{Statement and inherited foundation}
For finite $A\subset\mathbb R$, write
$\FS(A)=\{\sum_{x\in X}x:X\subseteq A\}$, including the empty sum,
and $4A=A+A+A+A$, allowing repeated summands. Let $N$ be the maximum
cardinality under $\FS(A)\subseteq4A$.
\begin{theorem}
Every feasible set has at most 95 elements. Hence $71\le N\le95$.
\end{theorem}
The lower bound is the unchanged construction
$[-31,31]\cap\Z\cup\{\pm93,\pm155,\pm217,\pm713\}$.
The upper-bound foundation is the independently audited chain
$115\to103\to100\to99\to98\to97\to96$; see \cite{previous,method} and
their dependencies. We reuse its integralization, fixed covering,
peeling, inverse, effective-hole, span, mass, and pairing lemmas.
We reprove here the additional geometric inequalities drawn from the
archived research report \cite{archive}, and generate fresh certificates.
No archived endpoint certificate is read.

A size-96 counterexample contains zero: otherwise adjoining zero gives
a forbidden size-97 example. This uses the preceding global bound;
feasibility is not assumed hereditary under arbitrary deletion.
Integralize, remove four outer elements on each side to obtain $D$,
and peel two largest and two smallest elements twice, obtaining
$D\supset D_1\supset E=P\cup\{0\}\cup(-Q)$. Put
\begin{gather*}
 d=88,\quad e=80,\quad m=a+b=79,\quad a=|P|,\quad b=|Q|,\\
 I=(a-39)(a-40),\quad T_r=r(r+1)/2,\\
 \Lambda=1+T_a+T_b=1601+I,\qquad \Omega=22d+97-\Lambda=432-I.
\end{gather*}

\section{A complete lattice gate at cardinality 96}
For a positive $r$-element set $X$, define
$\mu(X)=|\FS(X)|-1-T_r$ and
$\rho(X)=\mu(X)+2-2\delta_0(X)$, where $\delta_0(X)$ counts the $x$
whose double is absent from the distinct-element two-sum set of
$X\cup\{0\}$. Write $\mu=\mu(P)+\mu(Q)$ and
$\rho=\rho(P)+\rho(Q)$. With the inherited nonnegative quantities,
\begin{equation}\label{budget}
 \rho+2\vartheta+c+\T=\Omega,\quad
 c=2|2D|+18d+85-|\FS(D)|,\quad
 \T=\kappa+\ell_1+\ell_2+g.
\end{equation}
Here $\kappa=|\FS(E)|-\Lambda-\mu$; the $\ell_i$ are peeling slacks and
$g=4d-10-(|R(D)|-|R(D_1)|)$.
The deficient-sign case is impossible since
$1+T_{80}>24\cdot96+77$. The paired bootstrap gives $q^2\le2353$,
$q$ odd, hence $q\le47$ and $a,b\ge16$. Nonnegative $\Omega$ leaves
exactly $a=19,\ldots,60$. For these 42 oriented pairs,
\[
 T_{\min(a,b)}-2-\Omega\ge99>0.
\]
The simultaneous inverse theorem therefore gives a common core lattice,
with side-support holes satisfying
\[
 h_P,h_Q\le\kappa,\quad z=h_P+h_Q-\kappa,\quad
 0\le z\le\min(h_P,h_Q)\le\kappa,\quad h_X<|\FS(X)|-1.
\]
It remains to put the peeled elements on this lattice. The universal
span bound $\max X\le r+\rho(X)+h_X+2$, used before asserting lattice
membership of the peeled layers, gives, for $s=|\FS(E)|$ and
$C=2m+4e+14-\Lambda=-1109-I$,
\begin{align*}
 |\FS(D_1)|-2s&\le C+\Omega+\rho-\mu+2z,\\
 |\FS(D_1)|-2s&\le C+3\Omega-\rho-\mu.
\end{align*}
The maximum is always defective, so $\rho\le\mu$. For $z\le14$ the
first bound is at most $-649-2I<0$. For $z\ge15$ both supports are
non-progressions, so $\rho+\mu\ge(79-28)+(79-6)=124$.
The second bound is then $63-4I$. Any remaining pair has
$I=k(k+1)$ with $k\le3$, hence $a,b\ge36$.
The effective-hole row $(r,C_0,\beta)=(36,109,7)$ gives
$\mu(P)\ge a+110$, $\mu(Q)\ge b+110$; its strict tests are checked
exactly and remain valid for larger $r$. Thus $\rho+\mu\ge350$,
whereas $C+3\Omega=187-4I$. Across all pairs the resulting high-hole
comparison is at most $-17$.
Finally
$\Lambda-(4d-10+\Omega)=827+2I>0$.
The inherited overlap argument puts $D$ on the core lattice and
propagates its support holes. Only now do we use integer interval counts.

\section{A sharp four-extreme defect inequality}\label{sharp}
\begin{theorem}\label{sharp:thm:sharp}
Let $F=\{x_1<\cdots<x_m\}\subset\Z$, $m\ge7$, and let
\[
 a_1<a_2<x_1<\cdots<x_m<b_2<b_1,\qquad
 D=F\cup\{a_1,a_2,b_2,b_1\}.
\]
Let $\eta_F,\eta_D$ count the holes in the endpoint hulls of $R(F),R(D)$,
and put $g_H=4m+6-(|R(D)|-|R(F)|)$. Then
\begin{equation}\label{sharp:eq:sharp}
 \boxed{\eta_D+\eta_F+g_H\ge2m-2.}
\end{equation}
The constant is sharp for every $m\ge7$.
\end{theorem}
\begin{proof}
Write $L=x_1+x_2$ and $U=x_{m-1}+x_m$. Let $W_+$ be
$((b_1+F)\cup(b_2+F))\cap(U,\infty)$, and let $W_-$ be the corresponding
union of the two lower translates below $L$. Put
\[
 I_+=(U,b_1+x_m]\cap\Z,\quad I_-=[a_1+x_1,L)\cap\Z,
 \quad H_\pm=|I_\pm|-|W_\pm|.
\]
Let $q$ count the distinct mixed sums $a_i+b_j$ that lie outside $[L,U]$
and are not already in $W_-\cup W_+$. These sums lie in $I_-\cup I_+$,
so $0\le q\le4$ and $q\le H_-+H_+$.

Translate $x_1$ to zero and write
\[
\begin{gathered}
 x_m=M,\quad x_2=s_0,\quad x_{m-1}=M-t_0,\\
 a_2=-w_0,\quad a_1=-w_0-v_-,\quad
 b_2=M+u,\quad b_1=M+u+v_+.
\end{gathered}
\]
All six gap parameters are positive integers. At most $u+1$ members of
$b_2+F$ lie above $U$: at most $u$ eligible points are at or below
$x_{m-1}$, followed by $x_m$. Thus the integers
\[
 S_+=u+1+m-|W_+|,\qquad S_-=w_0+1+m-|W_-|
\]
are nonnegative. Set $H=H_++H_-$ and $S=S_++S_-$. We claim
\begin{equation}\label{sharp:eq:HS}
 H+S\ge2q-2.
\end{equation}
For $q\le2$ this follows from $H\ge q$. The same bound proves the case
$q=3,S\ge1$ and the case $q=4,S\ge2$.

Suppose first that $S=0$. Since $|W_\pm|\le2m$, we have
$u,w_0\le m-1$. The mixed value $B_0=a_2+b_2=M+u-w_0$ satisfies
\[
 M-m+2\le B_0\le M+m-2.
\]
Now $s_0\le M-m+2$ and $x_{m-1}\ge m-2$, so
$B_0\in[s_0,2M-t_0]=[L,U]$. In particular $q=4,S=0$ is impossible.

If $q=3,S=0$, the other mixed values
\[
 A_0=B_0-v_-,\qquad C_0=B_0+v_+-v_-,\qquad D_0=B_0+v_+
\]
are all outside $[L,U]$. Thus $A_0<L$ and $D_0>U$. Reflection allows us to
assume $C_0<L$. The definitions give
\[
 H_+=v_++t_0-m-1,\qquad H_-=v_-+s_0-m-1.
\]
Using integrality in $D_0>U$ and $C_0<L$ yields, respectively,
\[
 u-w_0\ge M-H_+-m,\qquad
 u-w_0\le t_0-M+H_--H_+-1.
\]
Consequently
\[
 H_-\ge2M-m-t_0+1=M+x_{m-1}-m+1\ge M-1\ge m-2\ge5.
\]
This proves \eqref{sharp:eq:HS} in that case.

Finally, suppose $q=4,S=1$. Reflect if necessary so that
$S_+=1,S_-=0$. The cardinality bounds now give $u\le m$ and
$w_0\le m-1$, so $B_0\ge M-m+2\ge s_0$. Since $B_0$ lies outside the
old hull, $B_0>U$, whence
\[
 u-w_0>x_{m-1}\ge m-2,\qquad u-w_0\le m-1.
\]
Equality forces $u=m$, $w_0=1$, and $x_{m-1}=m-2$. The first $m-1$
members of $F$ are therefore $0,1,\ldots,m-2$, so $s_0=1$, and
$B_0=M+m-1=U+1$. Because $v_-\ge1$, the value $A_0=B_0-v_-$ is at most
$U$; it must consequently be less than 1. Thus $v_-\ge M+m-1$ and
\[
 H_-=v_-+s_0-m-1=v_--m\ge M-1\ge m-2\ge5.
\]
Together with $S=1$, this proves \eqref{sharp:eq:HS} in the final case.

Put $h_-=L-(a_1+a_2)$, $h_+=(b_1+b_2)-U$, and $h=h_-+h_+$.
The definitions give the exact identity
\[
 2(|W_-|+|W_+|)=h+2m+2-H-S.
\]
Every new restricted sum outside $[L,U]$ lies in $W_-\cup W_+$, is one
of the $q$ new mixed values, or is one of $a_1+a_2,b_1+b_2$. The latter
values are distinct and outside $I_-\cup I_+$. Hence the number $N_{\rm out}$
of these sums satisfies
\[
 2N_{\rm out}=h+2m+6-H-S+2q\le h+2m+8.
\]
Write $N_0=|R(D)|-|R(F)|$. At most $\eta_F$ new sums are inside the old
hull, so $N_0\le N_{\rm out}+\eta_F$. Also
$\eta_D=h+\eta_F-N_0$. Therefore
\begin{align*}
 \eta_D+\eta_F+g_H
 &=h+2\eta_F+4m+6-2N_0\\
 &\ge h+4m+6-2N_{\rm out}\ge2m-2.
\end{align*}
For sharpness, take $F=[0,m-1]_{\Z}$ and
$\{a_1,a_2,b_2,b_1\}=\{-m-2,-1,m,2m+1\}$.
Here $R(F)=[1,2m-3]_{\Z}$ and $R(D)=[-m-3,3m+1]_{\Z}$.
For the latter assertion the four translates of $F$ fill the ranges
$[-m-2,-3]$, $[-1,m-2]$, $[m,2m-1]$, $[2m+1,3m]$;
$R(F)$ fills the central gap, and the sums of outer elements supply
$-m-3,-2,2m,3m+1$. Thus both hole counts vanish and $g_H=2m-2$.
\end{proof}

\section{Coupled budgets and endpoint transport}
Apply the sharp theorem to $F=D_1$, so
$G=2(d-4)-2=166$. The inherited hole propagation
$w\le4z$ yields
\begin{equation}\label{coupled}
 \T\ge\kappa+G-3z,\qquad
 \T\ge\max(56,z,166-2z).
\end{equation}
If $j$ of the eighteen covering units are empty, put $\delta=c-88j$.
Combining the span lemma with \eqref{budget}--\eqref{coupled} gives
\begin{equation}\label{span}
 R_p+R_q+\delta\le349-I-88j+4z,
 \qquad R_p=\max P,\quad R_q=\max Q.
\end{equation}
For $z=0$, $c\le266-I$; for $1\le z\le14$, both supports are
non-progressions and $\rho\ge51$, so $c\le215-I+2z\le243-I$.
Consequently the complete low-hole root ranges are
\begin{center}
\begin{tabular}{lrrrrrrr}\toprule
Regime&$Z_0$&$Z_1$&$Z_2$&$Z_3$&$P_0$&$P_1$&$P_2$\\\midrule
$a$&24--55&27--52&30--49&38--41&24--55&28--51&32--47\\
Integer roots&128&1820&11560&11968&128&1680&9248\\\bottomrule
\end{tabular}
\end{center}
These ranges deliberately retain every parameter allowed by the scalar
cap. In particular, $Z_1$ allows $\delta\le178\ge2\cdot84$ and $Z_2$
allows $\delta\le90\ge84$. Neither ``one large loss'' nor universal
$Z_2$ core quantiles may be assumed. In $Z_3$, $\delta\le2$, so ordinary
quantiles are valid. In $P_2$, $\delta\le67$, so core quantiles are valid.

For $z\ge15$, the checker reconstructs the inherited effective-hole
catalogue. For each side it chooses the strongest available certified
$u=r+C_0+1$ with $2\beta<z$, and checks the strict effective-hole tests,
including the endpoint bridge. The flexible inverse gives
\[
 \rho_*(r,u)=\max\{r-14,u+2-2r,L_2(r,u),\ldots,L_{10}(r,u)\},
\]
where, with $v=\max(0,u-4\beta+2)=(\beta+1)q+s$,
\[
 L_\beta(r,u)=\min\{(\beta-1)(r-14),(\beta-1)q+\min(s,\beta-1)\}.
\]
Thus $c\le\Omega-\rho_*(a,u_p)-\rho_*(b,u_q)-\max(56,z,166-2z)$.
There are 2,840 surviving oriented scalar states and $c\le144<176$,
so $j\le1$. For $j=0$, first orient the actual first-pair contributor
positively. In that fixed orientation the active bound is one of
\begin{align*}
 F&\ge\tfrac12T_a+T_b+6a+8b+85+\tfrac12(u_p+z)+u_q+z,\\
 F&\ge\tfrac23T_a+T_b+6a+(22b+236)/3+\tfrac23(u_p+z)+u_q+z,
\end{align*}
while $F\le1710+4z$. A positive minimum excludes that oriented state;
choosing the stronger reflected orientation is invalid.
This leaves 477 states for $j=0$ and 764 for $j=1$.
For $H_1$, $\delta\le56<84$ permits core quantiles.

\begin{lemma}[Mass transport]
For fixed side sizes and common lower bounds $u_p,u_q$, let $Z$ be the
largest hole count in a group and $C_*$ the largest covering cap.
Replace $S_p,S_q$ by $S_p+Z-z,S_q+Z-z$, leave all other coordinates
unchanged, and replace $z$ by $Z$. Every necessary row used below
continues to hold with cap $C_*$.
\end{lemma}
\begin{proof}
Writing $t=Z-z\ge0$, the doubled mass row has change
$-2t+2t=0$, the exposed mass row has change $-t+t=0$, the
cardinality row has change $2t-4t=-2t$, and the span row has change
$-4t$. Upper quantiles and leading-slot upper rows only weaken.
All overlap rows are unchanged. The low cap $\delta-2z$ decreases
by $2t$; the high cap is replaced by its maximum. Nonnegativity is
preserved. The shifted masses are auxiliary coordinates, not a claim
about a new actual set. All shifts are integral.
\end{proof}
For low positive holes take $Z=14$. For high holes group by
$(a,j,u_p,u_q)$, regardless of which affine part of the scalar cap is
active. This yields 12 $H_0$ and 14 $H_1$ groups, hence
$12\cdot4+14\cdot70=1028$ integer models instead of
$477\cdot4+764\cdot70=55388$.
The checker compares every source row with its transported row on all
1,241 retained states. This is in addition to the general coefficient
argument above.

\section{One loss budget for overlapping covering slots}
Let $S_p=\sum(D\cap\mathbb R_{>0})$ and
$S_q=-\sum(D\cap\mathbb R_{<0})$, initially before transport.
Write the original outer magnitudes as
$p_1>p_2>p_3>p_4>M_p$ and their negative-side analogues $q_i$.
Put $c_p=p_1-p_2-p_3-p_4$, $\alpha_p=p_2-p_3-p_4$,
$\beta_p=-p_1+2p_2-p_3-p_4$, and reflect for $q$.
The twenty center labels are
\[
 v_{p,i}=c_p+p_{i+1}\ (0\le i\le3),\quad
 v_{p,i+4}=c_p-q_{i+1}\ (0\le i\le3),\quad
 v_{p,8}=\alpha_p,\quad v_{p,9}=\beta_p,
\]
and their reflected counterparts. Pair 1 is $p4,q4$. In branch A,
$\beta_p+\beta_q\le0$, pair 2 is $p9,q9$; in branch B it is $p5,q5$.
There are sixteen singleton units and two paired units, each of capacity
$d$. Enumerate every $j$-subset of missing units. If pair 1 is present,
orient its actual contributor positively. If pair 2 is present and
not potentially split, select either actual contributor. In branch A
with $p5$ or $q5$ missing, pair 2 can split: do not count an individual
loss for it, but select an actual member and retain
$\beta_s+M_s\ge1$. The numbers of modes for $j=0,1,2,3$ are
$4,70,578,2992$; both member choices are retained.

For each retained individual center let $l_v,u_v$ be its point counts
below and above the target integer hull, and $\lambda_v=l_v+u_v$.
Let $K=|2D|$ and let $L_p,L_q$ be the losses of the leading slots
$c_p+2D,c_q+2D$ from their actual positive or reflected targets.
The fixed covering gives
\begin{equation}\label{loss}
 L_p+L_q+\sum_v\lambda_v\le\delta.
\end{equation}
We strengthen this by replacing groups of slots by their actual union
deficits. This construction is valid even when a leading slot is empty.

\begin{proposition}[Ordinary and leading groups]
For $k$ retained same-sign individual slots with extreme centers
$v_{\min},v_{\max}$, their union deficit is at least
\[
 H_G+\max_v\lambda_v,\qquad
 H_G=ke-(R_p+R_q+v_{\max}-v_{\min}+1),
\]
and at least $\sum_v\lambda_v$.
For a leading slot and $k$ retained own-center slots $c_p+p_i+D$,
the union deficit is at least
\[
 H_*+\max(L_p,\lambda_1,\ldots,\lambda_k),\qquad
 H_*=(k+2)e-2-p_{\max}-R_p-2R_q,
\]
and at least $L_p+\sum\lambda_i$. Reflect on the other side.
\end{proposition}
\begin{proof}
The core pieces of the first group lie in an integer interval of length
$R_p+R_q+v_{\max}-v_{\min}+1$; the other points number at most
$k(d-e)$. Subtract this raw union bound from $kd$. Points outside the
target include each member's loss set, adding at least its largest
loss. The sum bound follows by bounding the union by the sum of the
individual target contributions.
For the second group, $c_p+2E$ and the $c_p+p_i+E$ lie in
$[c_p-2R_q,c_p+p_{\max}+R_p]_{\Z}$. Since $|2E|\ge2e-1$, the other
points number at most $K-(2e-1)+k(d-e)$. Subtraction from $K+kd$
gives $H_*$. The same outside-target and sum arguments finish the proof.
\end{proof}
The leading centers are integral with at most one empty unit by the
inherited contributor argument. With two or three empty units,
$\delta\le90<175\le K$, so both leading slots contribute and their
centers are integral. Thus the integer hull count in this proposition
is justified in every regime.

The model partitions the retained centers, on each side, into own labels
$0,1,2,3$, cross labels $4,5,6,7$, and ungrouped labels $8,9$.
Each own block also includes its leading slot; at most three missing
units ensure that at least one own label remains. Introduce one
nonnegative variable $O$ for each own block and for each cross block of
size at least two. Choose $O$ to be that block's actual union deficit.
Impose its sum-of-losses bound and its ordinary hull bound for each
member. For every nonempty subset of an own block impose the leading
bound for that subset, plus the losses of the unselected members.
All these rows are valid: a union deficit for disjoint groups of slot
labels is at least the sum of their separate deficits. Finally,
\[
 \sum_{\text{blocks}}O+\sum_{\text{ungrouped }v}\lambda_v\le\delta.
\]
No overlap gain is counted twice. The variables $O$ represent actual
whole-block deficits; alternative subset bounds are inequalities for
the same variable, not separate credits added to the budget.

\section{Remaining necessary rows and finite disjunctions}
Write $x_0=R_p<x_1<x_2<x_3<x_4=M_p$ and similarly $y_i$.
All coordinates and losses are nonnegative. Besides the cap and span
rows, the model contains the inherited mass rows, on each side:
\begin{align*}
 2M_p+6R_p-2S_p+2z&\le-2T_a-12-2u_p,\\
 M_p+4R_p-S_p&\le-T_{a-1}-6,\\
 x_1+x_2+x_3+M_p-S_p+z&\le-T_a-u_p,\\
 R_p+x_1+x_2+x_3+M_p-S_p&\le-T_{a-1}.
\end{align*}
Here $u_p=0$ for $Z$, $a-3$ for $P$, and the certified effective-hole
bound for $H$; transport changes $z$ to its endpoint and the mass
constants accordingly. Also $R_p\ge a$, $M_p\ge R_p+4$, all inner gaps
are at least one, and
\[
 S_p+S_q+\delta-2K-4z\le1668-88j.
\]
Global positive-center quantiles are
$v+l_v+M_p\ge88$ and $v-u_v-M_q\le S_p-87$.
Core quantiles replace $M_p,M_q$ by $R_p,R_q$ and 88,87 by 84,83.
Ordinary quantiles are $v+l_v\ge b+5$, $v-u_v\le S_p-a-4$.
Use global quantiles generally, core quantiles for $P_2,H_1$, and
ordinary quantiles for $Z_3$. Individual loss caps are 87 globally and
83 in core models. Reflection supplies the other side.

If $j\ge2$, or the regime is $P_1$ or $H$, then $\delta<167$.
Counting the leading-slot contributors yields
\begin{align*}
 c_p+2M_p+L_p&\ge K,&S_p+2M_q-c_p+L_p&\ge K-1,&S_p+L_p&\ge K,\\
 c_p+2x_i+L_p&\ge167+2i,&S_p+2y_i-c_p+L_p&\ge166+2i.
\end{align*}
Indeed, a set consisting of $E$, four negative outliers and $i$ positive
outliers has at least $2(e+4+i)-1=167+2i$ twofold sums; at most $L_p$
fail to contribute. The bound $L_p\le\delta<167$ ensures a nonempty
counted interval. The three $K$ rows follow by counting all $K-L_p$
contributors. The coarse model replaces $K$ here by its lower bound 175.

With two empty units the leading centers are integral. If an original
outer element were nonintegral, its own unpaired slot would be empty.
There cannot be exactly one such element on a side, since $c_p$ or
$c_q$ is integral. Exactly two would exhaust the missing units; the
unpaired cross labels force outer elements 3 and 4 integral, and in
branch A the remaining cross label forces element 2 integral, while in
branch B the unpaired $\alpha$ does so. The leading-center identity then
forces element 1 integral, a contradiction. Thus strict outer gaps and
$\beta_p+\beta_q\ge1$ in branch B are valid for $j\le2$. For $j=3$
only weak outer gaps and a weak branch-B inequality are used.

The coarse envelope is $K\le M_p+M_q+R_p+R_q+21$; substitute it into
the cardinality row to eliminate $K$. Only the remaining difficult
groups retain inner coordinates and the five nested envelopes
\[
 K\le2(x_t+y_t)+1+T_{2(4-t)}+\sum_{s\in\{p,q\}}\sum_{i=t+1}^4g_{s,t,i},
\]
where $0\le g_{p,t,i}\le x_i-x_{i-1}$ and
$g_{p,t,i}\le x_t+y_t+1$, and similarly for $q$.
These are the inherited interval-union envelope rows. We also impose
\[
 K\le M_p+M_q+x_i+y_h+1+T_{4-i}+T_{4-h}\quad(0\le i,h\le4).
\]
Every mixed-sign sum and every same-sign sum with a summand below its
cutoff lies in $[-M_q-y_h,M_p+x_i]$. The only exceptions are pairs
among the remaining outliers, at most the displayed triangular counts.
This proves all 25 additional envelopes.

The finite trees use only the following exhaustive alternatives. All
statements apply equally after mass transport, with the enlarged upper
endpoint used to select an upper-endpoint alternative.
\begin{enumerate}
\item \emph{Core or large loss:} either $l_v+u_v\ge84$, or
$l_v+u_v\le83$ and both core quantiles hold. Actual losses are integers;
in the second case more than four points contribute, giving the core
counting inequalities. No bound on the number of large centers is needed.
\item \emph{Ordinary or outside:} either $v+l_v\ge b+5$, or
$l_v\ge b+5$ and $v\le-1$. At the upper end either
$v-u_v\le S_p-a-4$, or $u_v\ge a+5$ and $v\ge S_p+1$.
For nonnegative $v$, count the $b+5$ nonpositive members; for negative
$v$ all are lost. At the upper endpoint use the $a+4$ positive members
and, outside the hull, zero. Integral centers justify the unit gaps.
Enlarging the endpoint decreases the actual upper loss, so the original
loss variables still dominate the necessary counts.
\item \emph{Inner endpoint:} for $1\le i\le4$, either $l_v\ge i$ or
$v\ge y_{5-i}+1$; either $u_v\ge i$ or $v\le S_p-x_{5-i}$.
Otherwise all of the corresponding $i$ extreme points are lost.
\item \emph{Both or sole pair contributor:} in a split pair with selected
actual member $s$, either the other member satisfies the necessary
condition $\beta_t+M_t\ge1$, or $s$ is the sole contributor. In the
latter case introduce its two hull losses $l_*,u_*$, impose their global
quantiles, and add them to the sum-of-individual-losses budget.
Only necessity is used in the first branch; positivity alone is never
asserted sufficient for contribution.
\end{enumerate}
All alternatives keep both children. Since transformed masses, centers
and actual losses are integral, the displayed integer gaps remain valid.
Tree variables themselves may be relaxed to real nonnegative values.
\section{Exact certificate coverage and reproducibility}
The verifier reconstructs the complete root set from the scalar gates
and mode grammar. A root uses either the coarse model or the full
nested model. It does not enumerate designated large-loss positions:
where necessary, the proved core-or-large binary split supplies that
information locally. In the final certificate bundle:
\begin{center}
\begin{tabular}{lr}\toprule
Quantity&Count\\\midrule
Structural root groups&5,324\\
Integer root systems&37,560\\
Groups using the coarse model&5,301\\
Groups using the nested model&23\\
Direct symbolic interval checks&5,350\\
Integer roots requiring a tree&141\\
Binary internal nodes / tree leaves&275 / 416\\
Maximum tree depth&5\\
Final integer leaf systems&37,835\\
Distinct integer certificate vectors&651\\\bottomrule
\end{tabular}
\end{center}
The 275 splits consist of 207 core-or-large, 43 ordinary-or-outside,
24 both-or-sole pair, and one inner-endpoint split. The count identity
$416=141+275$ is verified. The minimum normalized strict certificate
gap is $1/73$.

For each system $Ax\le b$, $x\ge0$, a certificate is an integer row
weight vector $w\ge0$ with $w^TA\ge0$ coefficientwise and $w^Tb<0$.
Then $0\le w^TAx\le w^Tb<0$, a contradiction. The normalized gap is
$-w^Tb/\sum w_i$. In a parameter interval the matrix is fixed and
the right-hand sides are directly constructed as quadratic polynomials
in the integer side size $a$. Their weighted maximum is checked at
both endpoints and the relevant integers around the vertex, using
exact rational arithmetic. No polynomial is inferred from sampling.
There are 13,020 such extremum-point checks.

Every integer parameter is covered exactly once by a direct interval
or a complete two-child tree. Unknown rows, nonpositive contradictions,
negative weights, missing parameters, missing roots, unknown splits,
repeated splits on a path, and missing children are rejected. The
15 negative controls also reject omission of a three-empty mode,
omission of a split pair's actual contributor, the invalid universal
$Z_2$ core shortcut, the invalid one-large-loss shortcut, arbitrary
reflection in the active estimate, and an unsupported claim about
cardinality 95. In particular, this program certifies $N\le95$, not
$N\le94$.

\paragraph{Files and commands.}
The verifier is \path{verification/upper95.py}.\par
Its certificate data are \path{data/upper95_certificates.json}.\par
Verification needs Python 3.10 or later and only its standard library:
\begin{verbatim}
python -I -S -B verification/upper95.py --output .work/upper95.json
python -I -S -B verification/upper95.py --exhaustive
\end{verbatim}
The exhaustive mode regenerates every numerical integer model, compares
every row with the symbolic formula, and checks every assigned
certificate independently of interval compression. All tree leaves are
always checked at their actual integer parameter. The optional generator
\path{experiments/generate_upper95_certificates.py} uses SciPy to search
for multipliers and useful disjunctions; floating-point solver status
is not accepted as proof. Every saved multiplier is rechecked exactly.

The actual verification runs and hashes are stored under
\path{results/2026-10-16/upper95/}. Standard and optimized Python
outputs agree byte for byte; exhaustive and interval results agree
apart from the mode label; regenerated certificates agree byte for
byte; the preceding upper-96 result is unchanged. The Windows entry
\path{verify_upper95_windows.bat} runs the seven finite continuations
from 115 down to 95. Program success validates the finite certificates;
the infinite mathematical reductions are supplied by the report and
its cited foundation, rather than by a proof assistant.

\paragraph{What the method improves.}
Three-empty modes are handled explicitly with their tiny residual
budget. All other modes use a common loss geometry, so increasing the
allowed number of large-loss centers does not require a new global
case grammar. Only 23 of the 5,324 structural groups need nested geometry; every other
group closes with the coarse model. The high-hole transport independently
removes about 98.1\% of the unmerged high-hole model instances.

At the next cardinality 95, the balanced coarse parameters would be
$d=87$, $e=79$, $\Omega=450$, $G=164$. The low positive-hole cap can
reach $450-50-164+28=264\ge3\cdot87$, so positive three-empty modes
must be addressed as well. No exclusion of that endpoint is claimed.

\begin{thebibliography}{9}
\bibitem{previous} Sum-of-Subset research project,
\emph{Transported hole bounds and a proof of $N\le96$},
\path{reports/upper96_from97_20261016.tex}, 16 October 2026.
\bibitem{method} Sum-of-Subset research project,
\emph{Proof-method optimization},
\path{reports/proof_method_20261016.tex}, 16 October 2026.
\bibitem{archive} Sum-of-Subset research project,
\emph{Complete upper-86 research report},
\path{reports/upper86_complete_20261005.tex}, 5 October 2026.
The sharp four-extreme, grouped-overlap, leading-slot and outer-integrality
arguments are reconstructed here; its endpoint certificates are not used.
\end{thebibliography}
\end{document}
